% Options for packages loaded elsewhere
\PassOptionsToPackage{unicode}{hyperref}
\PassOptionsToPackage{hyphens}{url}
\documentclass[
  ignorenonframetext,
]{beamer}
\newif\ifbibliography
\usepackage{pgfpages}
\setbeamertemplate{caption}[numbered]
\setbeamertemplate{caption label separator}{: }
\setbeamercolor{caption name}{fg=normal text.fg}
\beamertemplatenavigationsymbolsempty
% remove section numbering
\setbeamertemplate{part page}{
  \centering
  \begin{beamercolorbox}[sep=16pt,center]{part title}
    \usebeamerfont{part title}\insertpart\par
  \end{beamercolorbox}
}
\setbeamertemplate{section page}{
  \centering
  \begin{beamercolorbox}[sep=12pt,center]{section title}
    \usebeamerfont{section title}\insertsection\par
  \end{beamercolorbox}
}
\setbeamertemplate{subsection page}{
  \centering
  \begin{beamercolorbox}[sep=8pt,center]{subsection title}
    \usebeamerfont{subsection title}\insertsubsection\par
  \end{beamercolorbox}
}
% Prevent slide breaks in the middle of a paragraph
\widowpenalties 1 10000
\raggedbottom
\AtBeginPart{
  \frame{\partpage}
}
\AtBeginSection{
  \ifbibliography
  \else
    \frame{\sectionpage}
  \fi
}
\AtBeginSubsection{
  \frame{\subsectionpage}
}
\usepackage{iftex}
\ifPDFTeX
  \usepackage[T1]{fontenc}
  \usepackage[utf8]{inputenc}
  \usepackage{textcomp} % provide euro and other symbols
\else % if luatex or xetex
  \usepackage{unicode-math} % this also loads fontspec
  \defaultfontfeatures{Scale=MatchLowercase}
  \defaultfontfeatures[\rmfamily]{Ligatures=TeX,Scale=1}
\fi
\usepackage{lmodern}
\ifPDFTeX\else
  % xetex/luatex font selection
\fi
% Use upquote if available, for straight quotes in verbatim environments
\IfFileExists{upquote.sty}{\usepackage{upquote}}{}
\IfFileExists{microtype.sty}{% use microtype if available
  \usepackage[]{microtype}
  \UseMicrotypeSet[protrusion]{basicmath} % disable protrusion for tt fonts
}{}
\makeatletter
\@ifundefined{KOMAClassName}{% if non-KOMA class
  \IfFileExists{parskip.sty}{%
    \usepackage{parskip}
  }{% else
    \setlength{\parindent}{0pt}
    \setlength{\parskip}{6pt plus 2pt minus 1pt}}
}{% if KOMA class
  \KOMAoptions{parskip=half}}
\makeatother
\setlength{\emergencystretch}{3em} % prevent overfull lines
\providecommand{\tightlist}{%
  \setlength{\itemsep}{0pt}\setlength{\parskip}{0pt}}
% Slide layout defaults for warning-clean scientific decks.
\usepackage{etoolbox}
\IfFileExists{xurl.sty}{\usepackage{xurl}}{}
\IfFileExists{seqsplit.sty}{\usepackage{seqsplit}}{\newcommand{\seqsplit}[1]{#1}}
\protected\def\breakseq#1{\seqsplit{#1}}
\protected\def\breaktt#1{\begingroup\ttfamily\seqsplit{#1}\endgroup}
\setlength{\emergencystretch}{6em}
\tolerance=5000
\hbadness=10000
\hfuzz=1pt
\setlength{\tabcolsep}{2pt}
\AtBeginEnvironment{longtable}{\tiny\renewcommand{\arraystretch}{0.86}\setlength{\tabcolsep}{1pt}}
\AtBeginEnvironment{tabular}{\tiny\renewcommand{\arraystretch}{0.86}\setlength{\tabcolsep}{1pt}}

% Natbib and cross-reference fallbacks — slides don't load natbib
% or cleveref, but combined-PDF manuscript prose may emit these
% commands. The fallback renders citations as a bracketed key list
% and cross-references as detokenized labels so slides don't fail on
% undefined control sequences or raw underscores. \providecommand is
% a no-op if packages load later.
\providecommand{\citep}[1]{[#1]}
\providecommand{\citet}[1]{#1}
\providecommand{\citealp}[1]{#1}
\providecommand{\citeauthor}[1]{#1}
\providecommand{\citeyear}[1]{#1}
\providecommand{\cref}[1]{\texttt{\detokenize{#1}}}
\providecommand{\Cref}[1]{\texttt{\detokenize{#1}}}
\usepackage{bookmark}
\IfFileExists{xurl.sty}{\usepackage{xurl}}{} % add URL line breaks if available
\urlstyle{same}
% fallback for those not using the hyperref driver hyperxmp:
\makeatletter
\@ifundefined{xmpquote}{\newcommand{\xmpquote}[1]{#1}}{}
\makeatother
\hypersetup{
  hidelinks,
  pdfcreator={LaTeX via pandoc}}

\author{\texorpdfstring{}{}}
\date{}

\begin{document}

\section{Results}\label{sec:results}

\begin{frame}[fragile,allowframebreaks]{Results}
\textbf{Run snapshot.} With the bundled \breaktt{manuscript/config.yaml}
the most recent execution evaluated the query \emph{``reproducible
research optimization''} against local, returned 6 deduplicated paper(s)
(4 carrying a DOI, 6 carrying an abstract); the per-source breakdown is
local=6 and recorded backend errors are none. The deep-search workflow
(\breakseq{[@sec:deep\_search]}) covered 3 keyword(s) --- \emph{convex
optimization; stochastic gradient descent; reproducible research} ---
drawn from arxiv, crossref, paperclip, producing unique paper(s) after
cross-keyword deduplication.

The diagnostic figures generated for this run are catalogued in
\breakseq{[@sec:methodology]}: \breakseq{[@fig:papers\_per\_source]} surfaces
per-backend coverage, \breakseq{[@fig:year\_histogram]} surfaces the temporal
distribution, and \breakseq{[@fig:score\_distribution]} surfaces the
relevance-score profile. The full determinism contract for each stage is
itemised in \breakseq{[@tbl:determinism]} of \breakseq{[@sec:reproducibility]}.
\end{frame}

\begin{frame}[fragile,allowframebreaks]{Interpreting the run snapshot}
\protect\phantomsection\label{interpreting-the-run-snapshot}
The numerical values in the run-snapshot paragraph that opens this
section are read directly from \breaktt{output/run\_summary.json} and
\breaktt{output/data/manuscript\_variables.json} so they always reflect
the most recent run rather than a stale claim hand-typed into the prose.
Three properties are worth highlighting (the formal determinism contract
for each underlying stage is enumerated in \breakseq{[@tbl:determinism]}):

\begin{itemize}
\tightlist
\item
  \textbf{Cache reuse is observable.} A second invocation against the
  same config produces a byte-identical artifact tree (modulo the
  wall-clock timestamp inside
  \texttt{output/search/cache/search\_\textless{}hash\textgreater{}.json}
  itself); see also the \emph{Search (cached hit)} row of
  \breakseq{[@tbl:determinism]} and the verification recipe in
  \breakseq{[@sec:reproducibility]}. The cache file thus doubles as a
  cryptographic seal: re-running is a file read, not a network
  round-trip.
\item
  \textbf{Deduplication is signal-preserving.} The aggregator merges by
  DOI → arXiv id → normalised (title, year), keeping the highest-scored
  copy and filling missing fields from the loser (see \emph{Dedup /
  merge} in \breakseq{[@tbl:determinism]} and the per-backend pre-dedup view in
  \breakseq{[@fig:papers\_per\_source]}). The \breaktt{RESULT\_NUM\_PAPERS}
  figure therefore equals ``papers a reviewer needs to read'', not ``raw
  backend hit count'' --- the per-source contributions in
  \breaktt{RESULT\_PER\_SOURCE} are the pre-dedup view.
\item
  \textbf{Enrichment coverage is honest.}
  \breaktt{RESULT\_WITH\_ABSTRACT} and \texttt{RESULT\_WITH\_DOI} count
  fields the corpus or the \texttt{AbstractFetcher} actually populated,
  never values inferred. When a paper is missing a DOI it is excluded
  from the DOI count even if its arXiv id resolves to one upstream. The
  temporal coverage of those papers is summarised by
  \breakseq{[@fig:year\_histogram]}, and their backend-reported relevance
  scores by \breakseq{[@fig:score\_distribution]}.
\end{itemize}
\end{frame}

\begin{frame}[fragile,allowframebreaks]{Output artefacts}
\protect\phantomsection\label{output-artefacts}
After running \breaktt{scripts/run\_search\_pipeline.py} against the
default \breaktt{manuscript/config.yaml}, the project produces:

\begin{itemize}
\tightlist
\item
  \breaktt{output/search/results.json} --- the raw \texttt{SearchResult}
  JSON, including \breaktt{per\_source\_counts} and \texttt{errors} for
  diagnostic purposes.
\item
  \texttt{output/search/cache/search\_\textless{}hash\textgreater{}.json}
  --- the deterministic search cache; identical reruns are file reads.
\item
  \texttt{output/cache/abs/\textless{}safe\_id\textgreater{}.txt} ---
  one file per fetched abstract.
\item
  \texttt{output/cache/pdf/\textless{}safe\_id\textgreater{}.\{pdf,txt\}}
  --- PDFs and extracted text (only when
  \breaktt{enrichment.fetch\_fulltext:\ true}).
\item
  \breaktt{output/corpus.json} --- a \texttt{LocalBackend}-compatible
  JSON corpus of every result, enriched in place.
\item
  \breaktt{manuscript/references.bib} --- the auto-populated bibliography
  from the single-query pipeline (merged with any other
  \breaktt{manuscript/*.bib} at PDF render time).
\item
  \texttt{output/llm/per\_paper/\textless{}safe\_id\textgreater{}.md}
  --- per-paper LLM analyses (only when \breaktt{llm.per\_paper:\ true}
  \emph{and} the LLM stack is reachable).
\item
  \breaktt{output/llm/synthesis.md} --- corpus-level LLM synthesis (only
  when \breaktt{llm.corpus\_synthesis:\ true} \emph{and} the LLM stack is
  reachable).
\item
  \breaktt{output/reading\_report.md} --- the final assembled reading
  report.
\end{itemize}

When the LLM stack is genuinely unreachable, the \texttt{output/llm/}
artefacts are simply absent --- no placeholder file is ever written into
the archive (see \breakseq{[@sec:pipeline\_internals]}).

Because the search cache and abstract cache are deterministic, a second
run with identical \texttt{config.yaml} produces byte-identical
artifacts (modulo timestamp metadata in the cache files themselves).
This is the property the project exists to demonstrate.

The exact paper count, DOI list, and synthesis text depend on the live
state of arXiv and Crossref at the time of the run --- and are therefore
not reproducible \emph{across} runs in different weeks. Users seeking
strict reproducibility should:

\begin{enumerate}
\tightlist
\item
  Pin a \texttt{LocalBackend} corpus generated from a successful run
  (\breaktt{infrastructure.search.literature.write\_corpus}) and remove
  \texttt{arxiv} / \texttt{crossref} from
  \breaktt{config.search.sources}.
\item
  Commit the \breaktt{output/search/cache/} directory to version control.
\item
  Pin the LLM seed (\texttt{config.llm.seed}) and avoid model upgrades.
\end{enumerate}

With those three steps, every run from the same commit produces the same
outputs.
\end{frame}

\end{document}
